\documentclass[10pt]{extarticle}
\usepackage[a4paper,margin=1.5cm]{geometry}
\usepackage[T1]{fontenc}
\usepackage{amsmath,amssymb}
\usepackage{graphicx}
\usepackage[export]{adjustbox}
\usepackage{multicol}
\usepackage{enumitem}
\usepackage{booktabs,array,tabularx}
\setlength{\columnsep}{0.7cm}
\setlength{\parindent}{0pt}
\setlength{\parskip}{3pt}
\pagestyle{empty}
\makeatletter
% figure registry (figdefs.tex is generated by the extract step)
\newcommand{\qfigdef}[4]{\@namedef{qf@file@#1}{#2}\@namedef{qf@w@#1}{#3}\@namedef{qf@wide@#1}{#4}}
\newcommand{\qf@check}[1]{\@ifundefined{qf@file@#1}{\PackageError{score2}{Unknown figure id #1}{}}{}}
\newcommand{\qf@img}[4][]{\includegraphics[width=#2\dimexpr\@nameuse{qf@w@#4}pt\relax,max width=#3,max height=0.45\textheight,#1]{\@nameuse{qf@file@#4}}}
% Full-width questions: finalize wraps every question that contains a wide figure in \qfullbegin...\qfullend,
% so the WHOLE question (stem, figures, options) is set across the page in reading order, then the two columns
% resume. The question's head - its number, stem and first figure (\qfullsplit marks the end of the head) -
% is kept on one page, so a number is never left at a page foot with its figure overleaf; the rest may break
% across pages normally, so no big gaps are left. Without a \qfullsplit, the whole question is kept together
% if it is under 45% of a page. Never write these yourself.
\newif\ifqf@full
\newif\ifqf@boxed
\newsavebox{\qf@box}
\newcommand{\qf@open}{\global\qf@fulltrue\par\global\qf@boxedtrue
  \setbox\qf@box\vbox\bgroup\hsize\textwidth\linewidth\textwidth}
\newcommand{\qf@emit}[1]{\par\egroup\global\qf@boxedfalse
  \ifdim\dimexpr\ht\qf@box+\dp\qf@box\relax<#1\textheight\noindent\usebox{\qf@box}\par\else\unvbox\qf@box\fi}
\newcommand{\qfullsplit}{\ifqf@boxed\qf@emit{0.85}\fi}
\newcommand{\qf@close}{\ifqf@boxed\qf@emit{0.45}\fi\global\qf@fullfalse}
\newcommand{\qfullbegin}{\end{multicols}\qf@open}
\newcommand{\qfullend}{\qf@close\begin{multicols}{2}}
\newcommand{\qfullstart}{\qf@open}
\newcommand{\qfulldone}{\qf@close}
% \qfig[scale]{id}: block figure, centred, directly where it is written (size set by the tool).
% Fallback if a wide figure is not inside a full-width question: break out of the columns in place - no floats.
\newcommand{\qfig}[2][1]{\qf@check{#2}%
  \ifqf@full\qf@blk{#1}{#2}\else
  \ifnum\@nameuse{qf@wide@#2}=1 \expandafter\@firstoftwo\else\expandafter\@secondoftwo\fi
  {\end{multicols}\begin{center}\qf@img{#1}{\textwidth}{#2}\end{center}\begin{multicols}{2}}%
  {\qf@blk{#1}{#2}}\fi}
\newcommand{\qf@blk}[2]{\par\smallskip{\centering\qf@img{#1}{\linewidth}{#2}\par}\smallskip}
% \qfigi[scale]{id}: inline figure (inside option lists / tables / minipages); never breaks columns.
\newcommand{\qfigi}[2][1]{\qf@check{#2}\qf@img[valign=c]{#1}{\linewidth}{#2}}
\makeatother
% structure macros - lint reads these to check completeness and figure placement
\newcommand{\papersec}[1]{\par\bigskip{\centering\Large\bfseries PAPER-#1\par}\medskip}
\newcommand{\partsec}[2]{\par\medskip\noindent{\large\bfseries PART-#1: #2}\par\smallskip}
\newcommand{\sect}[1]{\par\smallskip\noindent\textbf{#1}\par}
\newcommand{\qn}[1]{\par\smallskip\noindent\textbf{#1.}~}
\newcommand{\paraq}[2]{\par\smallskip\noindent\textbf{Paragraph for Questions #1 and #2}\par}
\newenvironment{qoptions}{\begin{itemize}[leftmargin=1.8em,itemsep=1pt,topsep=2pt,label={}]}{\end{itemize}}
\newcommand{\opt}[2]{\item[(#1)] #2}
% figure registry (generated by score2_tools.py finalize)
\qfigdef{p1_phy_s1_q01_1}{Score Advanced - Test 10 (06 May 2023)-img/p1_phy_s1_q01_1.png}{108.3}{0}
\qfigdef{p1_phy_s1_q03_1}{Score Advanced - Test 10 (06 May 2023)-img/p1_phy_s1_q03_1.png}{134.9}{0}
\qfigdef{p1_phy_s1_q05_1}{Score Advanced - Test 10 (06 May 2023)-img/p1_phy_s1_q05_1.png}{243.0}{0}
\qfigdef{p1_phy_s1_q06_1}{Score Advanced - Test 10 (06 May 2023)-img/p1_phy_s1_q06_1.png}{108.3}{0}
\qfigdef{p1_phy_s2_q07_1}{Score Advanced - Test 10 (06 May 2023)-img/p1_phy_s2_q07_1.png}{108.0}{0}
\qfigdef{p1_phy_s2_q08_1}{Score Advanced - Test 10 (06 May 2023)-img/p1_phy_s2_q08_1.png}{108.3}{0}
\qfigdef{p1_phy_s2_q09_1}{Score Advanced - Test 10 (06 May 2023)-img/p1_phy_s2_q09_1.png}{74.2}{0}
\qfigdef{p1_phy_s2_q10_1}{Score Advanced - Test 10 (06 May 2023)-img/p1_phy_s2_q10_1.png}{108.3}{0}
\qfigdef{p1_phy_s3_q01_1}{Score Advanced - Test 10 (06 May 2023)-img/p1_phy_s3_q01_1.png}{108.3}{0}
\qfigdef{p1_phy_s3_q03_1}{Score Advanced - Test 10 (06 May 2023)-img/p1_phy_s3_q03_1.png}{177.0}{0}
\qfigdef{p1_phy_s3_q04_1}{Score Advanced - Test 10 (06 May 2023)-img/p1_phy_s3_q04_1.png}{119.0}{0}
\qfigdef{p1_phy_s3_q05_1}{Score Advanced - Test 10 (06 May 2023)-img/p1_phy_s3_q05_1.png}{92.6}{0}
\qfigdef{p1_phy_s3_q06_1}{Score Advanced - Test 10 (06 May 2023)-img/p1_phy_s3_q06_1.png}{108.3}{0}
\qfigdef{p1_phy_s3_q07_1}{Score Advanced - Test 10 (06 May 2023)-img/p1_phy_s3_q07_1.png}{68.5}{0}
\qfigdef{p1_chem_s1_q01_1}{Score Advanced - Test 10 (06 May 2023)-img/p1_chem_s1_q01_1.png}{213.9}{0}
\qfigdef{p1_chem_s1_q02_1}{Score Advanced - Test 10 (06 May 2023)-img/p1_chem_s1_q02_1.png}{213.9}{0}
\qfigdef{p1_chem_s1_q02_2}{Score Advanced - Test 10 (06 May 2023)-img/p1_chem_s1_q02_2.png}{71.5}{0}
\qfigdef{p1_chem_s1_q02_3}{Score Advanced - Test 10 (06 May 2023)-img/p1_chem_s1_q02_3.png}{63.0}{0}
\qfigdef{p1_chem_s1_q04_1}{Score Advanced - Test 10 (06 May 2023)-img/p1_chem_s1_q04_1.png}{128.3}{0}
\qfigdef{p1_chem_s3_q01_1}{Score Advanced - Test 10 (06 May 2023)-img/p1_chem_s3_q01_1.png}{245.2}{0}
\qfigdef{p1_chem_s3_q01_2}{Score Advanced - Test 10 (06 May 2023)-img/p1_chem_s3_q01_2.png}{240.4}{0}
\qfigdef{p1_chem_s3_q03_1}{Score Advanced - Test 10 (06 May 2023)-img/p1_chem_s3_q03_1.png}{98.5}{0}
\qfigdef{p1_chem_s3_q03_2}{Score Advanced - Test 10 (06 May 2023)-img/p1_chem_s3_q03_2.png}{33.5}{0}
\qfigdef{p1_chem_s3_q04_1}{Score Advanced - Test 10 (06 May 2023)-img/p1_chem_s3_q04_1.png}{98.4}{0}
\qfigdef{p1_chem_s3_q04_2}{Score Advanced - Test 10 (06 May 2023)-img/p1_chem_s3_q04_2.png}{40.7}{0}
\qfigdef{p1_math_s3_q03_1}{Score Advanced - Test 10 (06 May 2023)-img/p1_math_s3_q03_1.png}{82.0}{0}
\qfigdef{p2_phy_s1_q01_1}{Score Advanced - Test 10 (06 May 2023)-img/p2_phy_s1_q01_1.png}{108.3}{0}
\qfigdef{p2_phy_s1_q04_1}{Score Advanced - Test 10 (06 May 2023)-img/p2_phy_s1_q04_1.png}{173.9}{0}
\qfigdef{p2_phy_s2_q06_1}{Score Advanced - Test 10 (06 May 2023)-img/p2_phy_s2_q06_1.png}{128.5}{0}
\qfigdef{p2_phy_s2_q07_1}{Score Advanced - Test 10 (06 May 2023)-img/p2_phy_s2_q07_1.png}{260.3}{1}
\qfigdef{p2_phy_s2_q08_1}{Score Advanced - Test 10 (06 May 2023)-img/p2_phy_s2_q08_1.png}{108.2}{0}
\qfigdef{p2_phy_s2_q09_1}{Score Advanced - Test 10 (06 May 2023)-img/p2_phy_s2_q09_1.png}{108.2}{0}
\qfigdef{p2_phy_s2_q10_1}{Score Advanced - Test 10 (06 May 2023)-img/p2_phy_s2_q10_1.png}{121.7}{0}
\qfigdef{p2_phy_s3_q02_1}{Score Advanced - Test 10 (06 May 2023)-img/p2_phy_s3_q02_1.png}{119.0}{0}
\qfigdef{p2_phy_s3_q03_1}{Score Advanced - Test 10 (06 May 2023)-img/p2_phy_s3_q03_1.png}{118.8}{0}
\qfigdef{p2_phy_s3_q04_1}{Score Advanced - Test 10 (06 May 2023)-img/p2_phy_s3_q04_1.png}{108.3}{0}
\qfigdef{p2_phy_s3_q05_1}{Score Advanced - Test 10 (06 May 2023)-img/p2_phy_s3_q05_1.png}{82.0}{0}
\qfigdef{p2_phy_s3_q06_1}{Score Advanced - Test 10 (06 May 2023)-img/p2_phy_s3_q06_1.png}{129.4}{0}
\qfigdef{p2_phy_s3_q08_1}{Score Advanced - Test 10 (06 May 2023)-img/p2_phy_s3_q08_1.png}{66.1}{0}
\qfigdef{p2_chem_s1_q04_A}{Score Advanced - Test 10 (06 May 2023)-img/p2_chem_s1_q04_A.png}{135.7}{0}
\qfigdef{p2_chem_s1_q04_B}{Score Advanced - Test 10 (06 May 2023)-img/p2_chem_s1_q04_B.png}{134.9}{0}
\qfigdef{p2_chem_s1_q04_C}{Score Advanced - Test 10 (06 May 2023)-img/p2_chem_s1_q04_C.png}{129.4}{0}
\qfigdef{p2_chem_s1_q04_D}{Score Advanced - Test 10 (06 May 2023)-img/p2_chem_s1_q04_D.png}{134.9}{0}
\qfigdef{p2_chem_s2_q09_A}{Score Advanced - Test 10 (06 May 2023)-img/p2_chem_s2_q09_A.png}{52.9}{0}
\qfigdef{p2_chem_s2_q09_2}{Score Advanced - Test 10 (06 May 2023)-img/p2_chem_s2_q09_2.png}{23.3}{0}
\qfigdef{p2_chem_s2_q10_1}{Score Advanced - Test 10 (06 May 2023)-img/p2_chem_s2_q10_1.png}{105.3}{0}
\qfigdef{p2_chem_s3_q07_1}{Score Advanced - Test 10 (06 May 2023)-img/p2_chem_s3_q07_1.png}{118.4}{0}
\qfigdef{p2_chem_s3_q07_2}{Score Advanced - Test 10 (06 May 2023)-img/p2_chem_s3_q07_2.png}{118.8}{0}
\qfigdef{p2_chem_s3_q07_3}{Score Advanced - Test 10 (06 May 2023)-img/p2_chem_s3_q07_3.png}{128.5}{0}
\qfigdef{p2_chem_s3_q07_4}{Score Advanced - Test 10 (06 May 2023)-img/p2_chem_s3_q07_4.png}{199.9}{0}
\qfigdef{p2_chem_s3_q08_1}{Score Advanced - Test 10 (06 May 2023)-img/p2_chem_s3_q08_1.png}{161.1}{0}
\qfigdef{p2_chem_s3_q08_2}{Score Advanced - Test 10 (06 May 2023)-img/p2_chem_s3_q08_2.png}{143.0}{0}
\raggedcolumns
\makeatletter\@beginparpenalty=10000\makeatother
\newcolumntype{L}[1]{>{\raggedright\arraybackslash}p{#1}}
\begin{document}
\begin{center}{\Large\bfseries Score Advanced --- Test 10 (06 May 2023)}\end{center}
\begin{multicols}{2}
\raggedcolumns
\papersec{1}
\partsec{1}{PHYSICS}
\sect{SECTION-I (i)}
\qn{1} Figure shows a potentiometer circuit. The length of potentiometer wire AB is 120 cm and its resistance is 2$\Omega$. The internal resistance of 2V cell is 1$\Omega$. Initially the switches S$_1$ and S$_2$ are open and S$_3$ is closed.\par
\qfig{p1_phy_s1_q01_1}
\begin{qoptions}
\opt{A}{When only S$_2$ is open, balanced length $\ell$ = 20 cm}
\opt{B}{When only S$_1$ is open, balanced length $\ell$ = 75 cm}
\opt{C}{When all three switches are closed, balanced length $\ell$ = 30 cm}
\opt{D}{When S$_3$ is open, S$_2$ is also open but S$_1$ is closed, balance length $\ell$ = 60 cm.}
\end{qoptions}
\qn{2} An $\alpha$-particle moves along a circle of radius 1.0 m in a homogeneous magnetic field of 1.0 Wb/m$^2$ :-\par
(h = 6.62 $\times$ 10$^{-34}$ J.s, e = 1.6 $\times$ 10$^{-19}$C, 1 amu = 1.6 $\times$ 10$^{-27}$ kg)
\begin{qoptions}
\opt{A}{De-Broglie wavelength for the particle will be 1.05 $\times$ 10$^{-16}$ m.}
\opt{B}{De-Broglie wavelength for the particle will be 2.06 $\times$ 10$^{-15}$ m.}
\opt{C}{Magnetic dipole moment of the particle will be 3.2 $\times$ 10$^{-11}$ Am$^2$ parallel to applied field.}
\opt{D}{Magnetic dipole moment of the particle will be 3.2 $\times$ 10$^{-11}$ Am$^2$ anti parallel to applied field.}
\end{qoptions}
\qn{3} Consider a long hypothetical cylindrical hollow planet of radius R and uniform mass per unit area $\sigma$. If a satellite revolves around the planet in circular orbit of radius r in a plane perpendicular to the axis of the planet, then :-\par
\qfig{p1_phy_s1_q03_1}
\begin{qoptions}
\opt{A}{kinetic energy of satellite is independent of orbital radius of satellite.}
\opt{B}{kinetic energy of satellite depends on radius of planet.}
\opt{C}{orbital speed of satellite depends on orbital radius of satellite.}
\opt{D}{time period revolution of satellite is directly proportional to its orbital radius.}
\end{qoptions}
\qn{4} Consider two sets of a measurement\par
Set-1 : 10.8 g, 10.7 g, 10.9 g\par
Set-2 : 10.431 g, 10.742g, 9.821 g\par
In set-1
\begin{qoptions}
\opt{A}{precision is greater in comparison to set-2.}
\opt{B}{precision is smaller in comparison to set-2.}
\opt{C}{Mean absolute error is greater comparison to set-2.}
\opt{D}{Mean absolute error is smaller comparison to set-2.}
\end{qoptions}
\qn{5} The common prism spectroscope is shown in figure is generally made up of three parts, viz., a prism and two telescopes. One of these telescopes is used to render the rays which fall on the prism parallel to each other, and is therefore called the collimator, the other telescope is made movable so that it can be turned into the proper direction for observing any desired color, and is, therefore, called the view-telescope. A circular scale is also provided to measure angular position of telescope. On the circular scale 29 divisions of the main scale exactly coincide with the 30 divisions of the vernier scale. The smallest division of the main scale is half-a-degree (0.5\textdegree). Initially view-telescope is focused to collect rays directly from collimator and position of telescope is measured. This is position (1). Now a prism is placed as shown in the figure and telescope is focused for red, yellow and violet. The respective positions of telescope are (2), (3) and (4).\par
{\footnotesize\setlength{\tabcolsep}{2pt}\noindent\begin{tabular}{|l|c|c|}
\hline
 & Reading of main scale & Reading of veriner scale\\ \hline
Position (1) & 20.0\textdegree & 20\\ \hline
Position (2) & 30.0\textdegree & 16\\ \hline
Position (3) & 32.5\textdegree & 9\\ \hline
Position (4) & 33.5\textdegree & 11\\ \hline
\end{tabular}}\par
\qfig{p1_phy_s1_q05_1}
The deviation for:
\begin{qoptions}
\opt{A}{red light 10\textdegree 4'}
\opt{B}{violet light 13\textdegree 19'}
\opt{C}{yellow light 12\textdegree 19'}
\opt{D}{violet light is 13\textdegree 21'}
\end{qoptions}
\qn{6} A solid non uniform sphere of mass M and radius R is released on smooth horizontal surface from the situation shown in the figure. In given situation centre of mass (C) and geometric centre (O) are on same horizontal line. The distance OC is $\dfrac{\text{R}}{2}$. The moment of inertia about axis passing through C and perpendicular to the plane of the paper is I$_0$ = $\dfrac{\text{MR}^2}{2}$. Acceleration of centre of mass is a and angular acceleration of the sphere is $\alpha$.\par
\qfig{p1_phy_s1_q06_1}
\begin{qoptions}
\opt{A}{Just after release a = $\sqrt{5}\dfrac{\alpha\text{R}}{2}$.}
\opt{B}{Just after release a = $\dfrac{\alpha\text{R}}{2}$.}
\opt{C}{At the moment when centre of mass is at the lowest point, angular velocity of the sphere is $\sqrt{\dfrac{2\text{g}}{\text{R}}}$.}
\opt{D}{At the moment when centre of mass is at the lowest point, angular velocity of the sphere is $\sqrt{\dfrac{\text{g}}{2\text{R}}}$.}
\end{qoptions}
\sect{SECTION-I (ii)}
\qn{7} The figure here shows a snapshot of a taut string carrying a wave traveling in positive x direction. Also shown are four identical elements named p, q, r and s on the string. List-I lists name of the elements, while list-II lists different characteristics. Considering upwards as positive, match entries in list-I to all possible entries in list-II.\par
\qfig{p1_phy_s2_q07_1}
{\footnotesize\setlength{\tabcolsep}{2pt}\noindent\begin{tabular}{|l|c|l|L{4.5cm}|}
\hline
 & List-I & & \multicolumn{1}{c|}{List-II}\\ \hline
(I) & p & (P) & Acceleration is toward negative y-axis.\\ \hline
(II) & q & (Q) & Velocity is toward positive y-axis.\\ \hline
(III) & r & (R) & Kinetic Energy is maximum.\\ \hline
(IV) & s & (S) & P.E. is maximum.\\ \hline
 & & (T) & Velocity is toward positive x-axis.\\ \hline
\end{tabular}}\par
\begin{qoptions}
\opt{A}{I $\to$ P;II $\to$ P,Q;III $\to$ Q,R,S;IV $\to$ Q}
\opt{B}{I $\to$ Q,R,S;II $\to$ P,Q;III $\to$ Q;IV $\to$ P}
\opt{C}{I $\to$ P;II $\to$ P,Q;III $\to$ Q,R;IV $\to$ T}
\opt{D}{I $\to$ T;II $\to$ P;III $\to$ Q;IV $\to$ Q}
\end{qoptions}
\qn{8} Volume versus pressure curves for n moles of an ideal gas are given for four process as shown in figure. AB is an isothermal curve. Match the entries of List-I with the enteries of List-II.\par
\qfig{p1_phy_s2_q08_1}
{\footnotesize\setlength{\tabcolsep}{2pt}\noindent\begin{tabular}{|l|l|l|L{4.2cm}|}
\hline
 & List-I & & \multicolumn{1}{c|}{List-II}\\ \hline
(I) & For process 1 & (P) & Work done is zero.\\ \hline
(II) & For process 2 & (Q) & Work done is maximum.\\ \hline
(III) & For process 3 & (R) & Internal energy will increase.\\ \hline
(IV) & For process 4 & (S) & Internal energy will decrease.\\ \hline
 & & (T) & Internal energy remains constant.\\ \hline
\end{tabular}}\par
\begin{qoptions}
\opt{A}{I $\to$ P,R;II $\to$ R;III $\to$ S;IV $\to$ P,S}
\opt{B}{I $\to$ Q,R;II $\to$ R;III $\to$ S;IV $\to$ P,Q}
\opt{C}{I $\to$ Q,R;II $\to$ S;III $\to$ R;IV $\to$ P,S}
\opt{D}{I $\to$ Q,R;II $\to$ R;III $\to$ S;IV $\to$ P,S}
\end{qoptions}
\qn{9} Identical capacitors A and B, each plate area S and separation between the plates d, are connected in parallel. Charge Q is given to each capacitor. Whole arrangement is shown in figure.\par
\qfig{p1_phy_s2_q09_1}
Now Keeping one plate of capacitor B fixed another plate is moved towards the first plate with very small constant velocity V from t = 0 and keeping one plate of A fixed, another plate is moved away from the first plate with small constant velocity V from t = 0. Match the proper entries from list-I to list-II for above described system.\par
{\footnotesize\setlength{\tabcolsep}{1.5pt}\noindent\begin{tabular}{|l|L{3.6cm}|l|l|}
\hline
 & \multicolumn{1}{c|}{List-I} & & List-II\\ \hline
(I) & Charge on capacitor A as function of time & (P) & $\dfrac{(\text{d}-\text{V t})\text{Q}}{\text{d}}$\\ \hline
(II) & Charge on capacitor B as function of time & (Q) & $\dfrac{(\text{d}+\text{V t})\text{Q}}{\text{d}}$\\ \hline
(III) & Current in the circuit as function of time & (R) & decreases\\ \hline
(IV) & Ratio of electrostatic potential energy stored between the plates of capacitor A and between the plates of capacitor B as function of time. & (S) & $\dfrac{\text{QV}}{\text{d}}$\\ \hline
 & & (T) & increases\\ \hline
\end{tabular}}\par
\begin{qoptions}
\opt{A}{I $\to$ Q,T;II $\to$ R;III $\to$ S;IV $\to$ P,R}
\opt{B}{I $\to$ P,R;II $\to$ Q,T;III $\to$ Q;IV $\to$ P}
\opt{C}{I $\to$ Q,T;II $\to$ S;III $\to$ P,R;IV $\to$ Q}
\opt{D}{I $\to$ P,R;II $\to$ Q,T;III $\to$ S;IV $\to$ R}
\end{qoptions}
\qn{10} A U shaped tube of mass 2m is placed on a horizontal surface. Two identical spheres each of diameter d (just less than the inner diameter of tube) and mass m enter into the tube with velocity v as shown in the figure taking all collisions to be elastic and all surfaces smooth.\par
\qfig{p1_phy_s2_q10_1}
{\footnotesize\setlength{\tabcolsep}{1.5pt}\noindent\begin{tabular}{|l|L{4.2cm}|l|l|}
\hline
 & \multicolumn{1}{c|}{List-I} & & List-II\\ \hline
(I) & The speed of the tube, when spheres are just about to collide inside the tube. & (P) & v\\ \hline
(II) & The speed of spheres, when spheres are just about to collide. & (Q) & 0\\ \hline
(III) & The speed of the spheres, when they come out of the tube. & (R) & $\dfrac{\sqrt{3}}{2}$v\\ \hline
(IV) & The speed of the tube, when spheres come out. & (S) & v/2\\ \hline
 & & (T) & 2v\\ \hline
\end{tabular}}\par
\begin{qoptions}
\opt{A}{I $\to$ S;II $\to$ R;III $\to$ Q;IV $\to$ P}
\opt{B}{I $\to$ S;II $\to$ R;III $\to$ P;IV $\to$ Q}
\opt{C}{I $\to$ S;II $\to$ Q;III $\to$ R;IV $\to$ P}
\opt{D}{I $\to$ R;II $\to$ S;III $\to$ P;IV $\to$ Q}
\end{qoptions}
\sect{SECTION-II}
\qn{1} Two unequal charged particles A and B are suspended with the help of two light inextensible string of length 3m and 6m respectively. Their masses are 1.6 kg and m$_2$ respectively. The distance between two points of suspension is 5m. In equilibrium two strings make 37\textdegree{} and 53\textdegree{} angle from vertical. Find the value of m$_2$ in grams.\par
\qfig{p1_phy_s3_q01_1}
\qn{2} A particle of m = 2kg is moving in a force field $\vec{\text{F}}$ = $-$8(x$\hat{\text{i}}$ + y$\hat{\text{j}}$)N. At t = 0 the particle is at (3m, 0) and its velocity is 4 m/s$\hat{\text{j}}$. Find the area enclosed by the path followed by the particle in SI units.
\qn{3} Consider following two circuits of ideal inductors and resistors. Both switches are closed at t = 0 and at this moment current is zero in all branches. If I$_1$ and I$_2$ are current in inductor of self inductance 2L in steady state for two circuits respectively, then find the value of $\dfrac{3\text{I}_1}{\text{I}_2}$.\par
\qfig{p1_phy_s3_q03_1}
\qn{4} Figure shows a solid hemispherical glass of refractive index $\mu$ = $\dfrac{3}{2}$ and radius R = $\dfrac{30}{\pi}$cm. A parallel beam of light is incident on the plane surface at angle of incidence 6\textdegree. At what distance (in cm) from principal axis beam will get focused after refraction from both surfaces?\par
\qfig{p1_phy_s3_q04_1}
\qn{5} Two balls are thrown with the same speed 50 m/s from the top of a cliff. The angles of their initial velocitites are 37\textdegree{} above and below the horizontal, as shown in figure. How much farther (in m) along the ground does the top ball hit than the bottom ball?\par
\qfig{p1_phy_s3_q05_1}
\qn{6} The system in figure lies on a frictionless horizontal table. Masses of two particles are 2m and m. The two spring constants are nk (where n is a numerical factor) and k. The springs have the same relaxed length. Assuming that it is possible to set up initial conditions so that the masses oscillate back and forth with the two springs always having equal lengths at any given instant, what is n?\par
\qfig{p1_phy_s3_q06_1}
\qn{7} Consider a conducting medium having non-uniform electric field in it. The specific resistance of medium is 5$\Omega$ m. Consider an irregular imaginary surface in medium. If electric current passing through the surface is 2 A then electric flux (in SI units) linked with the surface will be :\par
\qfig{p1_phy_s3_q07_1}
\qn{8} A source of alternating EMF is connected to a series combination of L, C \& R. If maximum and minimum value of instantaneous power delivered to the LCR circuit are 3P and -P, then phase difference between (in degree) applied EMF and current through circuit is :
\partsec{2}{CHEMISTRY}
\sect{SECTION-I (i)}
\qn{1} Choose the CORRECT option(s) regarding the given scheme\par
\qfig{p1_chem_s1_q01_1}
\begin{qoptions}
\opt{A}{[T] can show optical isomerism}
\opt{B}{[T] contains three functional groups}
\opt{C}{[T] can decolourise the KMnO$_4$ colour}
\opt{D}{[T] up on reaction with NH$_2$OH(excess) gives four diastereomers}
\end{qoptions}
\qn{2} \qfig{p1_chem_s1_q02_1}
Choose the correct option :
\begin{qoptions}
\opt{A}{x is Ph $-$ C $\equiv$ CH}
\opt{B}{y is \qfigi{p1_chem_s1_q02_2}}
\opt{C}{z is ester}
\opt{D}{x is \qfigi{p1_chem_s1_q02_3}}
\end{qoptions}
\qn{3} (i) $\psi_r = 2\left(\dfrac{1}{a_0}\right)^{3/2} e^{-r/na_0}$\par
(ii) $\psi_r = \dfrac{1}{9\sqrt{3}}\left(\dfrac{1}{a_0}\right)^{3/2}(6 - 6\sigma + \sigma^2)e^{-\sigma/2}$\par
(iii) $\psi_r = \dfrac{1}{2\sqrt{2}}\left(\dfrac{1}{a_0}\right)^{3/2}\left(2 - \dfrac{2r}{na_0}\right)e^{-r/na_0}$\par
(iv) $\psi_r = \dfrac{1}{9\sqrt{6}}\left(\dfrac{1}{a_0}\right)^{3/2}(4 - \sigma)\sigma e^{-\sigma/2}$, where $\sigma = \dfrac{2r}{na_0}$\par
Based on above radial functions for different orbitals, select correct statement(s) :
\begin{qoptions}
\opt{A}{for orbital (iv), the graph between $\psi_r$ and r will intersect x axis at two distinct non-zero finite values of r.}
\opt{B}{for orbital (i), the radial probability of finding electron will be maximum at r = a$_0$.}
\opt{C}{for orbital (ii), there is no angular node}
\opt{D}{for orbital (iii), distance of radial node from origin is `2a$_0$'}
\end{qoptions}
\qn{4} An ideal solution of two volatile liquids A and B is in equilibrium with its vapours as shown in figure. Let y$_1$ and x$_1$ be the mole fraction of A in vapour and liquid phase respectively. If the vapours are removed and condensed and equilibrium is re-established with remaining liquid solution. Let x$_2$ and y$_2$ be the mole fraction of A in liquid and vapour phase respectively at final equilibrium. Given : (P$^\text{o}_\text{A}$ $>$ P$^\text{o}_\text{B}$).\par
\qfig{p1_chem_s1_q04_1}
which of the following option(s) is/are always incorrect ?
\begin{qoptions}
\opt{A}{x$_2$ $>$ x$_1$}
\opt{B}{y$_2$ $<$ y$_1$}
\opt{C}{Vapour pressure of A will be more than that of B in reestablished final equilibrium condition.}
\opt{D}{Total pressure at the new equilibrium is less than that at initial equilibrium.}
\end{qoptions}
\qn{5} First species contains at least one weaker covalent bond than that present in second species in all of the following pairs EXCEPT?
\begin{qoptions}
\opt{A}{Diamond, graphite \ (C$-$C bond)}
\opt{B}{cyclopropane, propane \ (C$-$H bond)}
\opt{C}{HCO$_3^-$, CHO$_2^-$ \ (C$-$O bond)}
\opt{D}{PCl$_5$, PCl$_6^-$ \ (P$-$Cl bond)}
\end{qoptions}
\qn{6} Which statement(s) are true :
\begin{qoptions}
\opt{A}{In Bayer process, HCl can be used in place of CO$_2$.}
\opt{B}{In [Ni(CN)$_4$]$^{2-}$, d$_{\text{x}^2-\text{y}^2}$ of Ni$^{2+}$ participates in hybridisation.}
\opt{C}{Eu is largest lanthanoid.}
\opt{D}{[FeO$_4$]$^{2-}$ readily decomposes to give a paramagnetic gas.}
\end{qoptions}
\sect{SECTION-I (ii)}
\qn{7} Match the following List - I with List - II :\par
{\footnotesize\setlength{\tabcolsep}{1.5pt}\noindent\begin{tabular}{|c|L{3.9cm}|c|L{1.9cm}|}
\hline
 & \multicolumn{1}{c|}{List-I} & & \multicolumn{1}{c|}{List-II}\\ \hline
(P) & Neutralisation equivalent (E. Wt) is 59, on heating it gives another carboxylic acid of NE 74. & (1) & Maleic acid\\ \hline
(Q) & NE is 59, stable to heat, forms anhydride easily & (2) & Fumaric acid\\ \hline
(R) & NE is 58, stable to heat, forms anhydride easily & (3) & $\alpha$-methyl malonic acid\\ \hline
(S) & NE is 58, stable to heat, its pKa$_2$ is less than that of its diastereomer & (4) & Succinic acid\\ \hline
 & & (5) & Oxalic acid\\ \hline
\end{tabular}}\par
(Neutralisation equivalent)
\begin{qoptions}
\opt{A}{P $\longrightarrow$ 3;Q $\longrightarrow$ 4;R $\longrightarrow$ 2;S $\longrightarrow$ 1}
\opt{B}{P $\longrightarrow$ 4;Q $\longrightarrow$ 3;R $\longrightarrow$ 1;S $\longrightarrow$ 2}
\opt{C}{P $\longrightarrow$ 3;Q $\longrightarrow$ 4;R $\longrightarrow$ 1;S $\longrightarrow$ 2}
\opt{D}{P $\longrightarrow$ 4;Q $\longrightarrow$ 2;R $\longrightarrow$ 3;S $\longrightarrow$ 1}
\end{qoptions}
\qn{8} Match the list.\par
{\footnotesize\setlength{\tabcolsep}{1.5pt}\noindent\begin{tabular}{|c|L{3.6cm}|c|L{2.3cm}|}
\hline
 & \multicolumn{1}{c|}{List-I} & & \multicolumn{1}{c|}{List-II}\\ \hline
(P) & N$_2$O$_4$ $\rightleftharpoons$ 2NO$_2$ on increasing volume at constant temperature & (1) & Forward shift.\\ \hline
(Q) & H$_{2(g)}$ + I$_{2(g)}$ $\rightleftharpoons$ 2HI$_{(g)}$ on addition of inert gas at constant P \& T. & (2) & Concentration of reactant decreases.\\ \hline
(R) & 2NO$_{(g)}$ + Br$_{2(g)}$ $\rightleftharpoons$ 2NOBr$_{(g)}$ on decreasing pressure at constant temperature. & (3) & Concentration of product decreases.\\ \hline
(S) & 2ICl$_{3(g)}$ $\rightleftharpoons$ I$_{2(g)}$ + 3Cl$_{2(g)}$ [$\Delta$H $>$ 0] on increasing pressure at constant temperature. & (4) & Equilibrium is not disturbed.\\ \hline
 & & (5) & K$_\text{P}$ $>$ K$_\text{C}$ $\left(\text{T} > \dfrac{1}{\text{R}}\right)$\\ \hline
\end{tabular}}\par
\begin{qoptions}
\opt{A}{P $\longrightarrow$ 1,2,5;Q $\longrightarrow$ 2,3,4;R $\longrightarrow$ 2,3;S $\longrightarrow$ 3,5}
\opt{B}{P $\longrightarrow$ 1,2,5;Q $\longrightarrow$ 1,5;R $\longrightarrow$ 2,3;S $\longrightarrow$ 3,5}
\opt{C}{P $\longrightarrow$ 3;Q $\longrightarrow$ 1,4;R $\longrightarrow$ 1;S $\longrightarrow$ 2}
\opt{D}{P $\longrightarrow$ 1,2,3,5;Q $\longrightarrow$ 2,3,4;R $\longrightarrow$ 2,3;S $\longrightarrow$ 5}
\end{qoptions}
\qn{9} Match the list.\par
{\footnotesize\setlength{\tabcolsep}{1.5pt}\noindent\begin{tabular}{|c|L{3.6cm}|c|L{2.3cm}|}
\hline
 & \multicolumn{1}{c|}{List-I} & & \multicolumn{1}{c|}{List-II}\\ \hline
(P) & H$_{2(g)}$ at 0\textdegree{}C, 1 atm & (1) & Molar volume $<$ 22.4$\ell$\\ \hline
(Q) & Gas having density $>$ $\dfrac{10}{7}$g/$\ell$ at 0\textdegree{}C, 1 atm \& molecular weight of gas is 32. & (2) & Molar volume $>$ 22.4$\ell$\\ \hline
(R) & An unknown gas at 1 atm, 27\textdegree{}C (having Boyle temperature 0\textdegree{}C) & (3) & Behaves as an ideal gas.\\ \hline
(S) & A gas at very low density & (4) & Less compressible w.r.t. ideal gas\\ \hline
 & & (5) & a \& b are negligible.\\ \hline
\end{tabular}}\par
\begin{qoptions}
\opt{A}{P $\longrightarrow$ 4;Q $\longrightarrow$ 1,3,5;R $\longrightarrow$ 3,5;S $\longrightarrow$ 2,5}
\opt{B}{P $\longrightarrow$ 2,4;Q $\longrightarrow$ 1,4;R $\longrightarrow$ 2,3;S $\longrightarrow$ 5}
\opt{C}{P $\longrightarrow$ 2,4;Q $\longrightarrow$ 1;R $\longrightarrow$ 2,4;S $\longrightarrow$ 3,5}
\opt{D}{P $\longrightarrow$ 2,4;Q $\longrightarrow$ 1;R $\longrightarrow$ 3,5;S $\longrightarrow$ 4}
\end{qoptions}
\qn{10} List-I contains precipitate and List-II contains reagent in which it dissolves.\par
{\footnotesize\setlength{\tabcolsep}{1.5pt}\noindent\begin{tabular}{|c|l|c|l|}
\hline
 & List-I & & List-II\\ \hline
(P) & HgS & (1) & dil. HNO$_3$.\\ \hline
(Q) & Ni(OH)$_2$ & (2) & NaOH solution.\\ \hline
(R) & MnS & (3) & NH$_3$ solution.\\ \hline
(S) & Cr(OH)$_3$ & (4) & Na$_2$S (2M).\\ \hline
 & & (5) & Acetic acid solution.\\ \hline
\end{tabular}}\par
\begin{qoptions}
\opt{A}{P $\longrightarrow$ 1;Q $\longrightarrow$ 3;R $\longrightarrow$ 5;S $\longrightarrow$ 2}
\opt{B}{P $\longrightarrow$ 2;Q $\longrightarrow$ 1;R $\longrightarrow$ 3;S $\longrightarrow$ 4}
\opt{C}{P $\longrightarrow$ 4;Q $\longrightarrow$ 3;R $\longrightarrow$ 5;S $\longrightarrow$ 2}
\opt{D}{P $\longrightarrow$ 4;Q $\longrightarrow$ 2;R $\longrightarrow$ 1;S $\longrightarrow$ 5}
\end{qoptions}
\sect{SECTION-II}
\qn{1} An organic compound reacts with excess MeMgBr followed by H$_2$O to give two organic product of which one reacts with NaOI to give yellow precipitate and OTHER DOES NOT. How many of the following compounds can satisfy this condition ?\par
\qfig{p1_chem_s3_q01_1}
\qfig{p1_chem_s3_q01_2}
\qn{2} Organic compound [M] has the following data :\par
(i) Molecular weight is 149, it contains only 3 different elements.\par
(ii) It can rotate plane polarized light (chiral carbon).\par
(iii) [M] on oxidation with hot alkaline KMnO$_4$ produces benzoic acid.\par
(iv) [M] can be prepared by treating carbonyl compound with methyl amine followed by NaBH$_4$.\par
Write the number of possible structures for [M] (excluding stereo isomers).
\qn{3} What is the sum of position number assigned to Bromine atoms while numbering the parent chain in given structures :\par
\qfig{p1_chem_s3_q03_1}
\qfig{p1_chem_s3_q03_2}
\qn{4} If number of atoms in given unit cell whose centre lies on plane shown in figure A = x.\par
Number of ions whose centre lies on plane in a unit cell shown in figure B = y\par
Number of ions whose centres lie on plane in a unit cell shown in figure C = z\par
then (x + y + z) ?\par
\qfig{p1_chem_s3_q04_1}
\qfig{p1_chem_s3_q04_2}
\qn{5} To a solution of acetic acid, solid sodium acetate is added gradually forming a buffer. When x mole of salt is added, the pH has a certain value. If y mole of salt is added, the pH at this time changes by 0.6 units. Find the ratio (y/x). Given that y $>$ x.\par
[log2 = 0.3, log3 = 0.48, log5 = 0.7.]
\qn{6} When excess of KCN is added to aqueous copper sulphate solution, white ppt is obtained and after some time it disappear due to formation of `A'. How many electrons are present in d-orbitals of central metal ion of `A'.
\qn{7} How many of the following are redox reactions?\par
(i) Cu + dil. H$_2$SO$_4$ $\longrightarrow$\par
(ii) Zn + dil. H$_2$SO$_4$ $\longrightarrow$\par
(iii) H$_2$C$_2$O$_4$ + conc. H$_2$SO$_4$ $\longrightarrow$\par
(iv) P$_4$ + SOCl$_2$ $\longrightarrow$\par
(v) NH$_3$ + O$_2$ $\xrightarrow{\text{Pt}}$\par
(vi) Sn + PCl$_5$ $\longrightarrow$\par
(vii) H$_3$PO$_3$ $\xrightarrow{\Delta}$\par
(viii) Fe$^{3+}$ + SO$_2$ $\longrightarrow$\par
(ix) NaCl + H$_2$SO$_4$ $\longrightarrow$\par
(x) I$_2$ + H$_2$O($\ell$) $\longrightarrow$\par
(xi) HCl + O$_2$ $\xrightarrow{\text{CuCl}_2}$
\qn{8} PbS$_{(s)}$ + O$_{3(g)}$ $\longrightarrow$\par
Mass of Pb containing product in gram produced by each mole of O$_3$?\par
[Molecular weight : Pb = 207, S = 32, O = 16]
\partsec{3}{MATHEMATICS}
\sect{SECTION-I (i)}
\qn{1} Let a plane perpendicular to the plane 2x + 3y + 4z $-$ 5 = 0 contains the line $\dfrac{\text{x}-0}{5} = \dfrac{\text{y}-1}{3} = \dfrac{\text{z}-2}{4}$ and line L where line L is contained in the plane 2x + 3y + 4z $-$ 5 = 0, then :
\begin{qoptions}
\opt{A}{Line L passes through point $\left(0, \dfrac{7}{25}, \dfrac{26}{25}\right)$}
\opt{B}{Line L passes through point (0, 1, $-$1)}
\opt{C}{One of the direction ratios of line L is ($-$25, 6, 8)}
\opt{D}{One of the direction ratios of line L is ($-$25, $-$6, $-$8)}
\end{qoptions}
\qn{2} Which of the following is/are incorrect statement(s) ?
\begin{qoptions}
\opt{A}{If A is non-singular square matrix of order (n $\times$ n) and k is scalar (k, n) $\neq$ (0, 1) then $|$k adj(kA)$|$ is equal to k$^{n^2-n}|$A$|^{n-1}$.}
\opt{B}{Let A and B be 3 $\times$ 3 order matrices. If $|$A$|$ = 0, AB = 0 and AC = 0, then B = C.}
\opt{C}{A and B are two square matrices of order n $\times$ n such that AB = 0 and A is non-singular then B must be null matrix.}
\opt{D}{If AB is non-singular and A, B are square matrices then A and B both must be non-singular matrices.}
\end{qoptions}
\qn{3} If f(x) is a polynomial of degree three with leading coefficient as unity such that f(1) = 1, f(3) = 9 and f $'$(1) = 2, then :
\begin{qoptions}
\opt{A}{f $'$(x) = 2x has atleast one root between (1, 3).}
\opt{B}{f$''$(x) = 2 has atleast one root between (1, 7/3).}
\opt{C}{Product of roots of f(x) = 0 is 3.}
\opt{D}{Sum of roots of f(x) = 0 is 4.}
\end{qoptions}
\qn{4} A and B are the points (3, 0) and (0, 3) respectively. The coordinates of the point P on the line 3x + 4y + 1 = 0 is/are :
\begin{qoptions}
\opt{A}{(13, $-$10) if $|$PA $-$ PB$|$ is maximum.}
\opt{B}{$\left(-\dfrac{1}{7}, -\dfrac{1}{7}\right)$ if $|$PA $-$ PB$|$ is minimum.}
\opt{C}{(13, $-$10) if $|$PA $-$ PB$|$ is minimum.}
\opt{D}{$\left(-\dfrac{1}{7}, -\dfrac{1}{7}\right)$ if $|$PA $-$ PB$|$ is maximum.}
\end{qoptions}
\qn{5} Let f(x) = $\begin{cases} 5\text{x} & \text{x} \leq 2 \\ \displaystyle\int_0^{\text{x}} (4 + |2 - \text{t}|)\text{dt} & \text{x} > 2 \end{cases}$, then :
\begin{qoptions}
\opt{A}{f(x) is continuous but not differentiable at x = 2.}
\opt{B}{f(x) is not continuous at x = 2.}
\opt{C}{f(x) is differentiable everywhere $\forall$ x $\in$ (2, $\infty$)}
\opt{D}{The right hand derivative of f(x) at x = 3 does not exist.}
\end{qoptions}
\qn{6} Given that $f$(x) is a function defined on R, satisfying $f$(1) = 1 and for any x $\in$ R, $f$(x + 5) $\geq$ $f$(x) + 5 and $f$(x + 1) $\leq$ $f$(x) + 1. If g(x) = $f$(x) + 5 $-$ x, then g(2013) is greater than or equal to -
\begin{qoptions}
\opt{A}{3}
\opt{B}{4}
\opt{C}{5}
\opt{D}{6}
\end{qoptions}
\sect{SECTION-I (ii)}
\qn{7} Match the List :\par
{\footnotesize\setlength{\tabcolsep}{1.5pt}\noindent\begin{tabular}{|c|L{4.6cm}|c|c|}
\hline
 & \multicolumn{1}{c|}{List-I} & & \multicolumn{1}{c|}{List-II}\\ \hline
(I) & If $\displaystyle\int_a^b$ (f(x) $-$ 3x)dx = a$^2$ $-$ b$^2$ then value of f$\left(\dfrac{\pi}{6}\right)$ is : & (P) & $\dfrac{\pi}{3}$\\ \hline
(II) & The value of $\dfrac{\pi^2}{\ln 3}\displaystyle\int_{7/6}^{5/6}$ sec($\pi$x). dx is : & (Q) & $\dfrac{\pi}{6}$\\ \hline
(III) & If $\vec{\text{a}} = \hat{\text{j}} + \sqrt{3}\hat{\text{k}}$, $\vec{\text{b}} = -\hat{\text{j}} + \sqrt{3}\hat{\text{k}}$, $\vec{\text{c}} = 2\sqrt{3}\hat{\text{k}}$ form a triangle, then internal angle of the triangle between $\vec{\text{a}}$ and $\vec{\text{b}}$ is : & (R) & $\pi$\\ \hline
(IV) & $\displaystyle\int_0^\infty \dfrac{\text{dx}}{(1 + \text{x}^3)(1 + \text{x}^2)}$ = & (S) & $\dfrac{\pi}{4}$\\ \hline
 & & (T) & $\dfrac{2\pi}{3}$\\ \hline
\end{tabular}}\par
\begin{qoptions}
\opt{A}{I $\longrightarrow$ R;II $\longrightarrow$ S;III $\longrightarrow$ T;IV $\longrightarrow$ P}
\opt{B}{I $\longrightarrow$ Q;II $\longrightarrow$ R;III $\longrightarrow$ T;IV $\longrightarrow$ S}
\opt{C}{I $\longrightarrow$ P;II $\longrightarrow$ R;III $\longrightarrow$ S;IV $\longrightarrow$ T}
\opt{D}{I $\longrightarrow$ P;II $\longrightarrow$ R;III $\longrightarrow$ T;IV $\longrightarrow$ S}
\end{qoptions}
\qn{8} Match the list.\par
{\footnotesize\setlength{\tabcolsep}{1.5pt}\noindent\begin{tabular}{|c|L{5.2cm}|c|c|}
\hline
 & \multicolumn{1}{c|}{List-I} & & \multicolumn{1}{c|}{List-II}\\ \hline
(I) & If f(x) + f$''$(x) = $-$xg(x)f$'$(x) and g(x) $>$ 0, $\forall$ x $\in$ R then (f(x))$^2$ + (f$'$(x))$^2$ has its maxima at x = k, then (k$^2$ + 9) is : & (P) & 12\\ \hline
(II) & If f$^2$(x) + f(x) = g(x) and g(x) is always increasing, then the minimum value that f(x) can attain so that f(x) is also increasing is k. Then $|$8(k $-$ 1)$|$ is : & (Q) & 4\\ \hline
(III) & If g(x) = f(x) $\displaystyle\int_0^{\text{x}}$ f(t)dt is non-increasing function and f(x) is continuous in R, then (f(1.5) + 4) is : & (R) & 10\\ \hline
(IV) & If f(x) is differentiable function and f(x) = x$^2$ + $\displaystyle\int_0^{\text{x}}$ e$^{-\text{t}}$f(x $-$ t)dt, then f(2) = $\dfrac{\text{k}}{3}$, then k is : & (S) & 9\\ \hline
 & & (T) & 20\\ \hline
\end{tabular}}\par
\begin{qoptions}
\opt{A}{I $\longrightarrow$ R;II $\longrightarrow$ P;III $\longrightarrow$ T;IV $\longrightarrow$ Q}
\opt{B}{I $\longrightarrow$ S;II $\longrightarrow$ Q;III $\longrightarrow$ P;IV $\longrightarrow$ T}
\opt{C}{I $\longrightarrow$ S;II $\longrightarrow$ P;III $\longrightarrow$ Q;IV $\longrightarrow$ T}
\opt{D}{I $\longrightarrow$ S;II $\longrightarrow$ P;III $\longrightarrow$ R;IV $\longrightarrow$ T}
\end{qoptions}
\qn{9} Match the list.\par
{\footnotesize\setlength{\tabcolsep}{1.5pt}\noindent\begin{tabular}{|c|L{4.9cm}|c|c|}
\hline
 & \multicolumn{1}{c|}{List-I} & & \multicolumn{1}{c|}{List-II}\\ \hline
(I) & If the mid-point of chord of ellipse $\dfrac{\text{x}^2}{25} + \dfrac{\text{y}^2}{36} = 1$ is (0, 2), then the length of the chord is : & (P) & $\dfrac{3\sqrt{7}}{\sqrt{2}}$\\ \hline
(II) & The line 2x + y = a touches the hyperbola x$^2$ $-$ 2y$^2$ = 9, if $|$a$|$ is equal to : & (Q) & $2\sqrt{3}$\\ \hline
(III) & If y = $-$2x + $\lambda$ is a normal to the parabola y$^2$ = 4x, then value of $\sqrt{\lambda}$ is : & (R) & $\dfrac{20\sqrt{2}}{3}$\\ \hline
(IV) & The product of length of the perpendicular from the two focii on any tangent to the hyperbola $\dfrac{\text{x}^2}{9} - \dfrac{\text{y}^2}{15} = 1$ is k, then $\sqrt{\text{k}}$ is : & (S) & $\sqrt{30}$\\ \hline
 & & (T) & $\sqrt{15}$\\ \hline
\end{tabular}}\par
\begin{qoptions}
\opt{A}{I $\longrightarrow$ S;II $\longrightarrow$ Q;III $\longrightarrow$ P;IV $\longrightarrow$ T}
\opt{B}{I $\longrightarrow$ R;II $\longrightarrow$ Q;III $\longrightarrow$ T;IV $\longrightarrow$ P}
\opt{C}{I $\longrightarrow$ P;II $\longrightarrow$ T;III $\longrightarrow$ Q;IV $\longrightarrow$ S}
\opt{D}{I $\longrightarrow$ R;II $\longrightarrow$ P;III $\longrightarrow$ Q;IV $\longrightarrow$ S}
\end{qoptions}
\qn{10} Match the list.\par
{\footnotesize\setlength{\tabcolsep}{1.5pt}\noindent\begin{tabular}{|c|L{5.2cm}|c|c|}
\hline
 & \multicolumn{1}{c|}{List-I} & & \multicolumn{1}{c|}{List-II}\\ \hline
(I) & The value of $\lambda$ such that the planes x $-$ 2y $-$ z = 10, 3x + 3y + $\sqrt{\lambda}$z = 8 and 2x $-$ y + 2z = k, are not able to make faces of tetrahedron. & (P) & 8\\ \hline
(II) & The remainder when 3$^{113}$ is divided by 13 is : & (Q) & 6\\ \hline
(III) & Number of rational terms in the expansion of $\left(\sqrt{2} + \sqrt[3]{3} + \sqrt[6]{5}\right)^{12}$ is : & (R) & 9\\ \hline
(IV) & Let $\vec{\text{r}}_1, \vec{\text{r}}_2, \ldots, \vec{\text{r}}_n$ be the sides of a regular polygon inscribed in a circle if $|\vec{\text{r}}_1 \times \vec{\text{r}}_2 + \vec{\text{r}}_2 \times \vec{\text{r}}_3 + \ldots + \vec{\text{r}}_n \times \vec{\text{r}}_1| = |\vec{\text{r}}_1.\vec{\text{r}}_2 + \vec{\text{r}}_2.\vec{\text{r}}_3 + \ldots + \vec{\text{r}}_n.\vec{\text{r}}_1|$, then smallest possible value of n is : & (S) & 10\\ \hline
 & & (T) & 81\\ \hline
\end{tabular}}\par
\begin{qoptions}
\opt{A}{I $\longrightarrow$ T;II $\longrightarrow$ R;III $\longrightarrow$ Q;IV $\longrightarrow$ P}
\opt{B}{I $\longrightarrow$ R;II $\longrightarrow$ Q;III $\longrightarrow$ P;IV $\longrightarrow$ S}
\opt{C}{I $\longrightarrow$ S;II $\longrightarrow$ R;III $\longrightarrow$ Q;IV $\longrightarrow$ P}
\opt{D}{I $\longrightarrow$ T;II $\longrightarrow$ R;III $\longrightarrow$ P;IV $\longrightarrow$ S}
\end{qoptions}
\sect{SECTION-II}
\qn{1} A line L in the Argand plane is called as the mean line for points $\alpha_1, \alpha_2, \alpha_3, \ldots, \alpha_n$, if there exist points $\beta_1, \beta_2, \beta_3, \ldots, \beta_n$ on line L such that $\displaystyle\sum_{r=1}^{n}(\alpha_r - \beta_r) = 0$. If the line L passes through the point $-5 + 6i$, then the slope of the mean line L for points $-2 + 3i$, $5 + 8i$, $-6 + 60i$, $5 - 581i$ is $\lambda$, then value of [$\lambda$] is : (Where [.] greatest integer function)
\qn{2} A student can take only unit steps and each step in North, East, West and South with equal probability. He wants to go from his home at (2, 3) to his coaching at (4, 5). If he reaches by taking less than 7 steps with probability P, then 32P is ?
\qn{3} Number of ways in which 8 letters of word KATAPPAA can be placed in squares of the figure so that no row remains empty is $\lambda$, then $\dfrac{3}{73}\lambda$ is :\par
\qfig{p1_math_s3_q03_1}
\qn{4} The straight line $\dfrac{\text{x}}{4} + \dfrac{\text{y}}{3} = 1$ intersects the ellipse $\dfrac{\text{x}^2}{16} + \dfrac{\text{y}^2}{9} = 1$ at two points A and B, there is a point P on this ellipse such that the area of $\Delta$PAB is equal to 6$\left(\sqrt{2} - 1\right)$. Then the number of such points P is \_\_\_\_\_.
\qn{5} If 2 $\leq$ [x] $<$ 5, $-$2 $\leq$ [y] $<$ 4, 3 $\leq$ [z] $<$ 5, then maximum value of determinant $\left|\begin{matrix} [\text{x}]+2 & [\text{y}] & [\text{z}] \\ [\text{x}] & [\text{y}]+2 & [\text{z}] \\ [\text{x}] & [\text{y}] & [\text{z}]+2 \end{matrix}\right|$ is :\par
(where [.] denotes greatest integer function)
\qn{6} If the sum of all the solutions of the equation\par
$\left(2\sec^2\text{x} - \cos\dfrac{16}{9}\text{x}\right)\sec\text{x}$\par $- \left(\cos\dfrac{\text{x}}{3} + 2\tan\text{x}.\sec\text{x}\right)\tan\text{x} = \sin\dfrac{\text{x}}{3}$; x $\in$ [$-$9$\pi$, 9$\pi$] is k$\pi$, then k is equal to :
\qn{7} The differential equation $\dfrac{\text{dy}}{\text{dx}} = \dfrac{(\text{x}+\text{y}-3)^2 - (\text{y}-5)^2}{(\text{x}+2)^2}$ determines a curve y = f(x) and f(0) = 1, then $|f(3)|$ is equal to :
\qn{8} If p, q, r and s are non-zero numbers such that `r' and `s' are roots of x$^2$ + px + q = 0 and p, q are roots of x$^2$ + rx + s = 0, then value of (3p + 5q + 11r + 15s) is ?
\raggedcolumns
\papersec{2}
\partsec{1}{PHYSICS}
\sect{SECTION-I (i)}
\qn{1} A solid cylinder of mass m and radius R is placed on layer of liquid of viscosity $\eta$. The height of liquid is h (h $\ll$ R). Now the cylinder is rotated about its vertical axis with constant angular velocity $\omega$. Then :-\par
\qfig{p2_phy_s1_q01_1}
\begin{qoptions}
\opt{A}{External power required to rotate the cylinder with constant angular velocity $\omega$ is zero.}
\opt{B}{External power required to rotate the cylinder with constant angular velocity $\omega$ is $\dfrac{\pi\omega^2\text{R}^4\eta}{\text{h}}$.}
\opt{C}{External power required to rotate the cylinder with constant angular velocity $\omega$ is $\dfrac{\pi\omega^2\text{R}^4\eta}{2\text{h}}$.}
\opt{D}{External torque required to rotate the cylinder with constant angular velocity $\omega$ is $\dfrac{\pi\omega\text{R}^4\eta}{\text{h}}$.}
\end{qoptions}
\qn{2} Three co-hereent sources kept along the same line produce intensity I$_0$ each at point P on this line. When S$_1$ and S$_2$ are switched on simultaneously, intensity at point P is 2I$_0$. When S$_2$ \& S$_3$ are switched on simultaneously, intensity at point P is 2I$_0$. Then :-
\begin{qoptions}
\opt{A}{When S$_1$ and S$_3$ are switched on simultaneously, intensity at P can be 2I$_0$.}
\opt{B}{When S$_1$ and S$_3$ are switched on simultaneously, intensity at P can be 9I$_0$.}
\opt{C}{When all three sources are switched on simultaneously, intensity at P can be I$_0$.}
\opt{D}{When all three sources are switched on simultaneously, intensity at P can be 3I$_0$.}
\end{qoptions}
\qn{3} The pitch of a screw gauge is 1 mm and its circular scale is divided into 100 divisions. When nothing is placed between its studs, the zero of the circular scale lies 6 divisions below the reference line. When a wire is placed between its studs, the main scale reading is 2 mm and 56$^\text{th}$ division of circular scale coincides with the reference line of the main scale. The length of wire is 8.68 cm. Find the curved surface area (in cm$^2$) of the wire in true significant figures.
\begin{qoptions}
\opt{A}{6.82}
\opt{B}{6.28}
\opt{C}{8.26}
\opt{D}{8.82}
\end{qoptions}
\qn{4} In the circuit as shown in the adjoining figure the switch S$_1$ is closed at t = 0. After steady state is reached, the switch S$_2$ is also closed for a long time. Then the total heat dissipated in the circuit after the switch `S$_2$' is closed will be :\par
\qfig{p2_phy_s1_q04_1}
\begin{qoptions}
\opt{A}{100 $\mu$J}
\opt{B}{50 $\mu$J}
\opt{C}{150 $\mu$J}
\opt{D}{20 $\mu$J}
\end{qoptions}
\sect{SECTION-I (ii)}
\qn{5} A spherical ball of radius R is uniformly charged. The surrounding space of ball is charged with volume charge density $\dfrac{\alpha}{\text{r}}$, where $\alpha$ is a constant and r is the distance from the centre of the ball. The value of charge q on ball for which the strength of electric field outside the ball is independent of r. (The relative permittivites of the ball and the surrounding space are assumed to be equal to unity)
\begin{qoptions}
\opt{A}{Charge on ball is q = $\pi\alpha$R$^2$}
\opt{B}{Charge on ball is q = 2$\pi\alpha$R$^2$}
\opt{C}{Electric field out-side the ball is E = $\dfrac{2\alpha}{\varepsilon_0}$}
\opt{D}{Electric field out-side the ball is E = $\dfrac{\alpha}{2\varepsilon_0}$}
\end{qoptions}
\qn{6} A uniform disc of mass m and radius R lies on a board of mass M(M $>$ m), which lies on a frictionless inclined plane of inclination $\theta$, as shown. String and pulley are ideal. Assume that disc rolls without slipping on the board and assume that two string segments are parallel to the plane. Which of the following is/are correct?\par
\qfig{p2_phy_s2_q06_1}
\begin{qoptions}
\opt{A}{Acceleration of board is $\dfrac{(\text{M}-\text{m})\text{g}\sin\theta}{\text{M}+2\text{m}}$}
\opt{B}{Acceleration of board is $\dfrac{(\text{M}-\text{m})\text{g}\sin\theta}{\text{M}+3\text{m}}$}
\opt{C}{Angular acceleration of disc is $\dfrac{2(\text{M}-\text{m})\text{g}\sin\theta}{(\text{M}+3\text{m})\text{R}}$}
\opt{D}{Angular acceleration of disc is $\dfrac{(\text{M}-\text{m})\text{g}\sin\theta}{(\text{M}+3\text{m})\text{R}}$}
\end{qoptions}
\qfullbegin
\qn{7} Two particles A and B start to move from origin along +ve x-axis. In the time interval t = 0 to t = 2 sec. the velocity -time graph of both particles is shown in below figures.\par
\qfig{p2_phy_s2_q07_1}
\qfullsplit
\begin{qoptions}
\opt{A}{Separation between A and B is maximum at t = 1 sec.}
\opt{B}{Separation between A and B is maximum at t = 1.5 sec.}
\opt{C}{Maximum Separation between A and B is 2m}
\opt{D}{Maximum Separation between A and B is 2.5m}
\end{qoptions}
\qfullend
\qn{8} The figure shows the variation of current i of a discharging capacitor connected in series with resistance 10$\Omega$.\par
\qfig{p2_phy_s2_q08_1}
\begin{qoptions}
\opt{A}{The initial potential difference across the capacitor is 100 volt.}
\opt{B}{The capacitance of the capacitor is $\dfrac{1}{10\ell\text{n}2}$F.}
\opt{C}{The total heat produced in the circuit will be $\dfrac{500}{\ell\text{n}2}$ Joules.}
\opt{D}{The thermal power in the resistor will decrease with a time constant $\dfrac{1}{3\ell\text{n}2}$ second.}
\end{qoptions}
\qn{9} A potentiometer wire of resistance per unit length r = r$_0\left(1 + \dfrac{\text{x}}{4\ell}\right)$ (where x is distance from left end of the potentiometer wire) is connected to four ideal batteries as shown in the figure. Resistance of connecting wires is negligible. A cell of e.m.f. 1 volt is balanced against potential drop on potentiometer wire, then null point can be obtained :\par
\qfig{p2_phy_s2_q09_1}
\begin{qoptions}
\opt{A}{Between point A and B.}
\opt{B}{At point B.}
\opt{C}{Between point D and E.}
\opt{D}{At point C}
\end{qoptions}
\qn{10} A conducting rod of mass `m' and length `$\ell$' is placed on a pair of smooth horizontal conducting parallel rails connected at one end by an inductor of inductance L and by a capacitor of capacitance `C' at the other end as shown in the figure. A uniform vertical magnetic field `B' exists in the region perpendicular to the plane of rails. The conductor starts from rest at t = 0 from the left end of the rails due to constant external horizontal force F = 2mg applied to it. Then choose the correct option(s). (Take CB$^2\ell^2$ = m)\par
\qfig{p2_phy_s2_q10_1}
\begin{qoptions}
\opt{A}{The maximum horizontal displacement of the rod is $\dfrac{8\text{mgL}}{\text{B}^2\ell^2}$.}
\opt{B}{The maximum horizontal displacement of the rod is $\dfrac{4\text{mgL}}{\text{B}^2\ell^2}$.}
\opt{C}{The angular frequency of the oscillations of the rod is $\dfrac{\text{B}\ell}{\sqrt{\text{mL}}}$.}
\opt{D}{The angular frequency of the oscillations of the rod is $\dfrac{\text{B}\ell}{\sqrt{2\text{mL}}}$.}
\end{qoptions}
\sect{SECTION-III}
\qn{1} A small particle of mass m moves in such a way that the potential energy U = $\dfrac{1}{2}$mb$^2$r$^2$, where b is a constant and r is the distance of the particle from the origin (Nucleus). Assume Bohr model of quantization of angular momentum and circular orbits is applicable. If the radius of n$^\text{th}$ allowed orbit is proportional to n$^{4/\text{k}}$, then value of K is:-
\qn{2} A light string attached to a bead passing over a smooth fixed pulley at a distance of d = 15 m from the vertical rod as shown in the figure. At the other end of the string a block of mass 1 kg is attached, lying over a smooth fixed inclined plane of inclination angle 37\textdegree{}. Mass of the bead is 1 kg and friction between bead and vertical rod is zero. The bead is held in level with the pulley and released. If the speed of the bead when the string between bead and pulley makes an angle 53\textdegree{} with the horizontal is $\sqrt{\dfrac{\alpha \times 1000}{41}}$m/s, find value of $\alpha$. [take g = 10 m/s$^2$]\par
\qfig{p2_phy_s3_q02_1}
\qn{3} In a region of free space, a uniform magnetic field of induction B varies with time t as shown in the graph. At the instant t = 0, a particle of charge q and mass 2 kg enters the region with velocity 4m/s perpendicular to the direction of the magnetic field, neglect the influence of the induced electric field and explore motion of the particle for what minimum value of charge q (in C), will the particle follow a closed path ?\par
\qfig{p2_phy_s3_q03_1}
\qn{4} Three small balls A, B and C of masses m, 2m and 2m respectively are placed on a smooth horizontal surface. Balls A and B as well as B and C are connected by light inextensible strings each of equal length $\ell$. The string connecting balls A and B is taut. The initial distance between the balls B and C is $\dfrac{\ell\sqrt{3}}{2}$ and BC is perpendicular to AB. The ball `C' is given a horizontal velocity v$_0$ = 23 m/s parallel to AB as shown in the figure. Find the initial velocity of ball A (in m/s) just after string BC becomes taut.\par
\qfig{p2_phy_s3_q04_1}
\qn{5} A satellite is projected into space from a point `P' at a distance `r$_0$' from the centre of the earth with an initial velocity v$_0$ = $\sqrt{\dfrac{\text{GM}}{\text{r}_0}}$ at an angle $\alpha$ = 37\textdegree{} with the horizontal at point `P' as shown in the figure. It is found that the satellite is propelled into an elliptical orbit. If the eccentricity of the elliptical orbit is e, then value of 5e is :-\par
\qfig{p2_phy_s3_q05_1}
\qn{6} Two wires of same length $\ell$ = 1 m and same cross-sectional area A = 10 mm$^2$ are welded together and their ends are rigidly clamped between two walls as shown in the figure. Given :\par
$\alpha_1$ = 2 $\times$ 10$^{-5}$ C$^{-1}$, Y$_1$ = 1 $\times$ 10$^{10}$ N/m$^2$, $\mu_1$ = 0.4 kg/m\par
$\alpha_2$ = 1 $\times$ 10$^{-5}$ C$^{-1}$, Y$_2$ = 2 $\times$ 10$^{10}$ N/m$^2$, $\mu_2$ = 0.1 kg/m\par
The tension in the wires is produced by reducing the temperature of system by 20\textdegree{}C. If the first overtone frequency of transverse vibration of the system with joint `O' as a node is 5n Hz, then find the value of n.\par
\qfig{p2_phy_s3_q06_1}
\qn{7} A metallic surface is illuminated alternatively with electromagnetic radiations of wavelengths 4500 \AA{} and 6000 \AA{}. It is observed that the maximum speeds of the photoelectrons emitted under these illuminations are in the ratio 2 :1. Find the work function (in eV in nearest integer) of the metal. (Take hC = 1242 eV-nm)
\qn{8} A uniform stick with mass m and length 6 meter lies on a frictionless horizontal surface (so you can ignore gravity in this problem). Stick is pivoted at a point a distance x from its center. A spring (at its relaxed length) with spring constant k is attached to the far end of the stick, perpendicular to the stick, as shown in figure. If the stick is given a tiny kick, what value of x (in meter) yields the largest frequency of small oscillations?\par
\qfig{p2_phy_s3_q08_1}
\raggedcolumns
\partsec{2}{CHEMISTRY}
\sect{SECTION-I (i)}
\qn{1} Which of the following is correct order of solubility of given salts in aqueous solution ?
\begin{qoptions}
\opt{A}{LiClO$_4$ $<$ NaClO$_4$ $<$ KClO$_4$ $<$ RbClO$_4$}
\opt{B}{CaCO$_3$ $<$ NaHCO$_3$ $<$ KHCO$_3$ $<$ RbHCO$_3$}
\opt{C}{LiF $>$ NaF $>$ KF $>$ RbF $>$ CsF}
\opt{D}{AgF $<$ AgCl $<$ AgBr $<$ AgI}
\end{qoptions}
\qn{2} For an electrolytic cell undergoing following balanced reaction :\par
NO$_3^-$ + 3H$_3$O$^+$ + Hg$_2^{2+}$ $\rightarrow$ 2Hg$^{2+}$ + HNO$_2$ + 4H$_2$O\par
Under standard state conditions E$^0_\text{cell}$ = 0.02V . If an average current of 10 amp. passes through the cell, how long will it take to produce 0.3mole of HNO$_2$ by this reaction?\par
[Given NO$_3^-$ + 3H$_3$O$^+$ + 2e$^-$ $\rightarrow$ HNO$_2$ + 4H$_2$O, E$^0$ = 0.94V and $\dfrac{2.303\text{RT}}{\text{F}}$ = 0.06]
\begin{qoptions}
\opt{A}{1.6 hours}
\opt{B}{16 hours}
\opt{C}{3.2 hours}
\opt{D}{32 hours}
\end{qoptions}
\qn{3} H$_3$C $-$ C $\equiv$ C $-$ OCH$_3$ $\xrightarrow{\text{Hg}^{2+}/\text{H}^+/\text{H}_2\text{O}}$ P, P is :
\begin{qoptions}
\opt{A}{H$_3$C $-$ C(=O) $-$ CH$_2$ $-$ OCH$_3$}
\opt{B}{H$_3$C $-$ CH$_2$ $-$ C(=O) $-$ OCH$_3$}
\opt{C}{H$_3$C $-$ C(=O) $-$ C(=O) $-$ OCH$_3$}
\opt{D}{H$_3$C $-$ CH(OH) $-$ CH$_2$ $-$ OCH$_3$}
\end{qoptions}
\qn{4} Monomer units of Biodegradable polymer PHBV are :
\begin{qoptions}
\opt{A}{\qfigi{p2_chem_s1_q04_A}}
\opt{B}{\qfigi{p2_chem_s1_q04_B}}
\opt{C}{\qfigi{p2_chem_s1_q04_C}}
\opt{D}{\qfigi{p2_chem_s1_q04_D}}
\end{qoptions}
\sect{SECTION-I (ii)}
\qn{5} Which of the following statement(s) is are correct ?
\begin{qoptions}
\opt{A}{2ClO$_2$ + 8H$^+$ + 10I$^-$ $\longrightarrow$ 5I$_2$ + 2Cl$^-$ + 4H$_2$O}
\opt{B}{ClO$_2$ is diamagnetic in nature}
\opt{C}{ClO$_2$ in alkaline solution undergoes disproportionation}
\opt{D}{Chlorine dioxide (ClO$_2$) is powerful oxidising agent but bleaching action is lower than Cl$_2$}
\end{qoptions}
\qn{6} Which of the following are correct statements(s) :
\begin{qoptions}
\opt{A}{CoCl$_3$.3NH$_3$ shows facial and meridional isomerism.}
\opt{B}{[Co(gly)$_3$] shows geometrical isomerism as well as optical isomerism.}
\opt{C}{K$_4$[Fe$^\text{II}$(CN)$_5$(O$_2$)] is paramagnetic.}
\opt{D}{[Co(EDTA)]$^-$ shows optical isomerism.}
\end{qoptions}
\qn{7} Which of the following statement(s) is/are False ?
\begin{qoptions}
\opt{A}{If heat of vaporization of water at 100\textdegree{}C is 40.6 KJ/mol, then when 9gm of water vapours are condensed to liquid at 100\textdegree{}C \& 1 atm, then $\Delta$S$_\text{system}$ = 54.42 J/K}
\opt{B}{When ($\Delta$G$_\text{system}$)$_\text{P,T}$ $<$ 0 ; the reaction must be exothermic.}
\opt{C}{dG = Vdp $-$ SdT is applicable for closed system in which both PV and non-PV work occurs}
\opt{D}{If ($\Delta$G$_\text{system}$)$_\text{P,T}$ $<$ 0, then process can be work producing}
\end{qoptions}
\qn{8} To an aqueous solution containing a mixture of 0.05M NaCl and 0.05M NaI when AgNO$_3$ is slowly added to such a solution, then correct option(s) is/are :\par
(K$_\text{sp}$AgCl = 10$^{-10}$, K$_\text{sp}$AgI = 4 $\times$ 10$^{-16}$)
\begin{qoptions}
\opt{A}{Concentration of Ag$^+$ required to precipitate Cl$^\ominus$ is 2 $\times$ 10$^{-9}$M}
\opt{B}{Concentration of Ag$^+$ required to precipitate I$^\ominus$ is 8 $\times$ 10$^{-15}$M}
\opt{C}{AgCl and AgI will precipitate out of the solution together}
\opt{D}{First AgI will precipitate}
\end{qoptions}
\qn{9} Which oxidasing agent can not be used to convert Propanal to Propanoic acid :
\begin{qoptions}
\opt{A}{\qfigi{p2_chem_s2_q09_A}}
\opt{B}{HNO$_3$}
\opt{C}{SeO$_2$}
\opt{D}{CrO$_3$ / \qfigi{p2_chem_s2_q09_2}}
\end{qoptions}
\qn{10} Which reagent can be used to form 2-Chloropropane from Propan-2-ol
\begin{qoptions}
\opt{A}{NaCl}
\opt{B}{(i) \qfigi{p2_chem_s2_q10_1}}
\opt{C}{NH$_4$Cl}
\opt{D}{PCl$_3$}
\end{qoptions}
\sect{SECTION-III}
\qn{1} Consider Al$_2$(OH)$_6$ compound and calculate the value of $\dfrac{(\text{X} + \text{Y})}{\text{Z}}$.\par
Where X = Total number of (2C $-$ 2e$^-$) bond\par
Y = Total number of (3C $-$ 2e$^-$) bond\par
Z = Total number of (3C $-$ 4e$^-$) bond
\qn{2} How many of following elements produces NO$_2$ gas on reaction with concentrated HNO$_3$?\par
Zn, Fe, Pb, Cu, Ag, Au, Pt, Sn
\qn{3} Find the number of impurities which deposited as anode mud in electrorefining of copper.\par
Sb, Se, Te, Ag, Pt, Al
\qn{4} An electron in a He$^+$ ion is in a state with directional dependence and has a total of one node. A photon of frequency $\dfrac{5\text{RC}}{9}$ causes this election to jump into a higher state which can have total five possible orientations in space. Find magnitude of the change in (n + $\ell$) values corresponding to two states ?\par
(R = Rydbergs constant, C = speed of light)
\qn{5} 1g of a monobasic acid dissolved in 200g of water lowers its freezing point by 0.186\textdegree{}C. On the other hand when 1g of same acid is dissolved in 200ml of water, this solution requires 125ml of 0.1N NaOH for complete neutralization. If degree of ionisation of this acid in water is `$\alpha$' find value of `10$\alpha$' ?\par
[Given : For water K$_\text{f}$ = 1.86 K-Kg mol$^{-1}$]
\qn{6} For the reaction : 2NO(g) + H$_2$(g) $\rightarrow$ N$_2$O(g) + H$_2$O(g), the value of $\dfrac{-\text{dp}}{\text{dt}}$ was found to be 1.5 PaS$^{-1}$ for a pressure of 372 Pa for NO and 0.25 PaS$^{-1}$ for a pressure of 152 Pa for NO, the pressure of H$_2$ being constant. If pressure of NO was kept costant, the value of $\dfrac{-\text{dp}}{\text{dt}}$ was found to be 1.60 PaS$^{-1}$ for a pressure of 289 Pa for H$_2$ and 0.79PaS$^{-1}$ for a pressure of 144 Pa for H$_2$. If order of reaction with respect to NO and H$_2$ are a and b, respectively, then find nearest integer value of (a + b) ?
\qn{7} Consider following reactions \& write number of reaction which are given with correct major product\par
\qfig{p2_chem_s3_q07_1}
\qfig{p2_chem_s3_q07_2}
\qfig{p2_chem_s3_q07_3}
\qfig{p2_chem_s3_q07_4}
\qn{8} Consider following compounds \& write no. of compounds which can show faster SN$_1$ reaction than\par
\qfig{p2_chem_s3_q08_1}
\qfig{p2_chem_s3_q08_2}
\raggedcolumns
\partsec{3}{MATHEMATICS}
\sect{SECTION-I (i)}
\qn{1} The value of log$_{10}\left(\dfrac{\displaystyle\sum_{i=0}^{r}{}^n\text{C}_{2i}\,{}^{n-2i}\text{C}_{r-i}}{\displaystyle\sum_{i=r}^{n}{}^n\text{C}_i\,{}^{2i}\text{C}_{2r}\left(\dfrac{3}{4}\right)^{n-i}\left(\dfrac{1}{4}\right)^{i-r}}\right)$ (n $\geq$ 2r) is :
\begin{qoptions}
\opt{A}{2}
\opt{B}{1}
\opt{C}{0}
\opt{D}{None}
\end{qoptions}
\qn{2} A Point S moves in space so that the points A, B, C taken on X, Y, Z axes respectively are such that SA, SB, SC are mutually perpendicular and volume of tetrahedron OABC is 6, where O is origin then locus of S is :
\begin{qoptions}
\opt{A}{$\left(\text{x}^2 + \text{y}^2 + \text{z}^2\right)^3$ = 6xyz}
\opt{B}{$\left(\text{x}^2 + \text{y}^2 + \text{z}^2\right)^3$ = 168xyz}
\opt{C}{$\left(\text{x}^2 + \text{y}^2 + \text{z}^2\right)^3$ = 288xyz}
\opt{D}{$\left(\text{x}^2 + \text{y}^2 + \text{z}^2\right)^3$ = 48xyz}
\end{qoptions}
\qn{3} The coefficient of x$^{45}$ in $\displaystyle\prod_{i=1}^{10}\left(\text{x}^i - \text{i}\right)$ is :
\begin{qoptions}
\opt{A}{$-$8}
\opt{B}{$-$6}
\opt{C}{6}
\opt{D}{2}
\end{qoptions}
\qn{4} Let $\alpha, \beta, \gamma$ be non-zero real roots of the equation x$^3$ + $\alpha$x$^2$ + $\beta$x + $\gamma$ = 0 then
\begin{qoptions}
\opt{A}{There are three such triplets ($\alpha, \beta, \gamma$)}
\opt{B}{There are two such triplets ($\alpha, \beta, \gamma$)}
\opt{C}{There are eight such triplets ($\alpha, \beta, \gamma$)}
\opt{D}{There are four such triplets ($\alpha, \beta, \gamma$)}
\end{qoptions}
\sect{SECTION-I (ii)}
\qn{5} Consider the function g(x) = $\dfrac{f^2(\text{x})}{2}$ + 5x$^5$ $-$ 5 where f(x) is continuous and differentiable function. It is given that g(x) is monotonically increasing function and f(0) = 6 then
\begin{qoptions}
\opt{A}{g(3) $>$ 13}
\opt{B}{g$\left(\dfrac{1}{2}\right)$ $>$ 11}
\opt{C}{f$^{\,2}$(1) $>$ 26}
\opt{D}{f$^{\,2}$($-$1) $<$ 46}
\end{qoptions}
\qn{6} If \par\resizebox{\linewidth}{!}{$\displaystyle\lim_{n\to\infty}\left(\dfrac{1 + \sqrt[3]{2} + \sqrt[3]{3} + \ldots + \sqrt[3]{n}}{n^{7/3}\left(\dfrac{1}{(\alpha n+1)^2} + \dfrac{1}{(\alpha n+2)^2} + \ldots + \dfrac{1}{n^2(\alpha+1)^2}\right)}\right)$}\par $=$ 15 and $\alpha \in \text{R}^+$ then $\alpha$ is equal to :
\begin{qoptions}
\opt{A}{characteristic value of log$_{10}$3$^9$}
\opt{B}{$\displaystyle\lim_{\text{x}\to 0}\left(\dfrac{\sin(4\sin\text{x})}{\text{x}}\right)$}
\opt{C}{number of non-differentiable points of f(x) = $\left|\text{x}^2 - 3\,|\text{x}| + 2\right|$}
\opt{D}{$\left|\text{Re}\left(\sqrt{23 - 10\sqrt{2}\text{i}}\right)\right|$ (where i = $\sqrt{-1}$)}
\end{qoptions}
\qn{7} If $\alpha_n = \displaystyle\int_0^{I/2}\cos^n\text{x}\cos\text{nx dx}$ and $\beta_n = \displaystyle\int_0^{1}\dfrac{\sin(\text{nx})}{\sin\text{x}}\text{dx}$ then for the sequence $\{\alpha_n\}$ and $\{\beta_n\}$ :\par
\{Where I = K$\pi$, K $\in$ N\}
\begin{qoptions}
\opt{A}{$\alpha_1$ , $\alpha_2$ ..., $\alpha_n$ are in GP}
\opt{B}{$\beta_1$ , $\beta_2$ ..., $\beta_n$ are in GP}
\opt{C}{If I = $\pi$ then $\displaystyle\sum_{\text{i}=1}^{2000}\beta_\text{i}$ = 1000$\pi$}
\opt{D}{If I = $\pi$ then $\alpha_{100}$ = $\dfrac{\pi}{2^{100}}$}
\end{qoptions}
\qn{8} Let f : R $\longrightarrow$ R is given by\par
f(x) = $\begin{cases} (\text{x}+1)^3 + \text{x}, & \text{x} < 0 \\ \text{x}^2 - \text{x} + 1, & 0 \leq \text{x} < 1 \\ \dfrac{1}{2}\text{x}^2 - 4\text{x} + \dfrac{9}{2}, & 1 \leq \text{x} < 3 \\ (\text{x}-2)\ell\text{n}(\text{x}-2) - \text{x}, & \text{x} \geq 3 \end{cases}$\par
then which of the following option is/are correct?
\begin{qoptions}
\opt{A}{f has a point of local maximum and a point of local minimum in ($-\infty$ ,0)}
\opt{B}{f is an onto function}
\opt{C}{f is an increasing on ($-\infty$ ,0)}
\opt{D}{f $'$ is not differentiable at x = 1}
\end{qoptions}
\qn{9} For the matrix A = $\resizebox{!}{2.6\baselineskip}{$\begin{bmatrix} -1 & 0 & 0 \\ -1 & 0 & -1 \\ 0 & -1 & 0 \end{bmatrix}$}$ which of the following is/are true.
\begin{qoptions}
\opt{A}{A$^{2\text{K}}$ = $\begin{bmatrix} 1 & 0 & 0 \\ \text{K} & 1 & 0 \\ \text{K} & 0 & 1 \end{bmatrix}$ , K $\in$ N}
\opt{B}{A$^{2\text{K}-1}$ = $\begin{bmatrix} -1 & 0 & 0 \\ -\text{K} & 0 & -1 \\ -\text{K}-1 & -1 & 0 \end{bmatrix}$ , K $\in$ N}
\opt{C}{A$^{100}$ = $\begin{bmatrix} 1 & 0 & 0 \\ 50 & 1 & 0 \\ 25 & 0 & 1 \end{bmatrix}$}
\opt{D}{A$^{99}$ = $\begin{bmatrix} -1 & 0 & 0 \\ -50 & 0 & -1 \\ -49 & -1 & 0 \end{bmatrix}$}
\end{qoptions}
\qn{10} Pushpa is standing at (1,2) (height of pushpa is negligible). Another person Konda reddy is moving on a circular path (x $-$ 6)$^2$ + (y $-$ 8)$^2$ = 4. Pushpa wants to kill the Konda reddy by firing the bullet than which of the following is/are correct :\par
(Assuming path of bullet is straight line).\par
(Possibility of killing is inversely proportional to distance between pushpa \& Konda)
\begin{qoptions}
\opt{A}{The minimum angle pushpa can fire with +ve x-axis is tan$^{-1}\left(\dfrac{15 + \sqrt{228}}{21}\right)$.}
\opt{B}{The maximum angle pushpa can fire with +ve x-axis is tan$^{-1}\left(\dfrac{30 + \sqrt{228}}{21}\right)$.}
\opt{C}{The minimum distance between pushpa and Konda is $\left(\sqrt{61} - 2\right)$.}
\opt{D}{The least possibility to kill the Konda is when he is at the point $\left(\dfrac{10}{\sqrt{61}} + 6\ ,\ \dfrac{12}{\sqrt{61}} + 8\right)$.}
\end{qoptions}
\sect{SECTION-III}
\qn{1} If $\displaystyle\int\dfrac{\left(\sqrt{t} - 1\right)\text{dt}}{\left(\text{t} + \text{t}^{5/4} + \text{t}^{3/4}\right)\sqrt{\text{t}^{1/4}\left(\text{t}^{1/2}+1\right)}}$ = 8tan$^{-1}$ (g(t)) + C and g(1) = $\sqrt{2}$ then the value of 2g$^2$(16) is :\par
(where C is integration constant)
\qn{2} A continuous y = f(x) is defined in a closed interval [$-$7, 6]. If A($-$7,$-$4), B($-$2, 6), C(0, 0), D(1, 6), E(2, $-$3), F(3, 6), G(5,$-$6) and H(6, 0) are consecutive points on the graph of f(x) and if these points are joined as line segments AB, BC, CD, DE, EF, FG, GH, then graph of f(x) is formed and if number of real solutions of f(f(x)) = 6 is $\lambda$ then $\lambda$ $-$ 10 is equal to :
\qn{3} Let g$_1$ , g$_2$ , g$_3$ : (0, $\infty$) $\longrightarrow$ R be defined by g$_1$(x) = $\displaystyle\int_0^{\text{x}}\prod_{\text{i}=1}^{19}(\text{t}-\text{i})^\text{i}\text{dt}$, x $>$ 0\par
and g$_2$(x) = 98(x $-$ 1)$^{50}$ $-$ 100(x $-$ 1)$^{49}$ + 2450, x $>$ 0\par
and g$_3$(x) = (x $-$ 1)$^{2020}$ (x $-$ 2)$^{2021}$ (x $-$ 3)$^{2022}$, x $>$ 0\par
Let $\alpha_i$ and $\beta_i$ respectively denote the number of points of local minima and the number of points of local maxima of the function g$_i$ , i = 1, 2, 3 in (0, $\infty$) and if the value of $\alpha_1\beta_1$ + $\alpha_2\beta_2$+ $\alpha_3\beta_3$ is $\lambda$ then $\lambda$ $-$ 20 equal to ( where $\displaystyle\prod_{i=1}^{n}$ a$_i$ = a$_1$a$_2$a$_3$ . . . . . a$_n$, n $\in$ N)
\qn{4} The area bounded by the curve f(x) = x + sinx and its inverse between the ordinates x = 0 to x = 3$\pi$ is 3K, K $\in$ N then K is
\qn{5} X wants to put 6 different rings in 3 specified fingers. The number of different ways he can put such that no finger remains empty and no finger can have more than 3 rings is 6! $\times$ $\lambda$. Then the value of $\lambda$ is : ($\lambda$ $\in$ N)
\qn{6} If y = g(x) is the solution of the differential equation x(y$^3$ $-$x)dy = y (x + y$^3$)dx and g(1) = ($-$2)$^{\frac{1}{3}}$ then $\dfrac{g^{-1}(-4)}{16}$ is equal to : (where g$^{-1}$ denotes the inverse of g(x))
\qn{7} The number of solution(s) of the equation $\left|\tan^{-1}2\,|\text{x}|\right|$ = $\sqrt{16\text{x}^4 - 8\text{x}^2 + 1}$, is(are)
\qn{8} Let f(x) be an invertible function and f(f$^{-1}$(x)) = f$^{-1}$ (f(x)) = x for f(x) = x$^3$ $-$ 2x$^{-3}$ + 10 (x $>$ 0), then the value of 9$\dfrac{\text{d}}{\text{dx}}$(f$^{-1}$(x)) at x = 9 is :
\end{multicols}
\end{document}
